\section{先认识你正在设计的东西}
\subsection{四个实验各增加了什么能力}
实验1接收两个数字，再根据开关选择结果。实验2让电路保存数据。实验3把保存的数据逐拍加一，用它衡量时间，再控制灯的位置。实验4将按键处理、计数和显示连在一起，形成一个能交互的小系统。

\begin{center}
\begin{tikzpicture}[node distance=5mm]
\node[block,minimum width=34mm] (a) {实验1：组合逻辑\\计算与选择};
\node[block,right=of a,minimum width=34mm] (b) {实验2：寄存器\\保存数据};
\node[block,below=9mm of b,minimum width=34mm] (c) {实验3：计数器\\计时与状态更新};
\node[block,left=of c,minimum width=34mm] (d) {实验4：控制器\\按键与多位显示};
\draw[flow] (a)--(b);\draw[flow] (b)--(c);\draw[flow] (c)--(d);
\end{tikzpicture}
\end{center}

\begin{center}
\begin{tabularx}{\linewidth}{@{}lXX@{}}\toprule
实验 & 要回答的问题 & 输出取决于什么\\\midrule
1 & 现在应该输出加法还是减法结果？ & 当前输入与控制开关\\
2 & 之前存了什么，现在允许改写吗？ & 保存的数据、写入条件和时钟\\
3 & 什么时候移动，下一步亮哪里？ & 计数进度、当前位置和控制开关\\
4 & 怎样把数显示出来，把一次按下识别出来？ & 多个模块各自的状态与连接\\\bottomrule
\end{tabularx}
\end{center}

\subsection{FPGA、开发板和Verilog分别是什么}
FPGA是一种可以配置内部逻辑功能与连接方式的芯片。EGO1是开发板的名称，板上除了FPGA，还有时钟源、开关、按键、LED和数码管。你的目标器件是\code{xc7a35tcsg324-1}。

Verilog是硬件描述语言。你用它描述输入输出关系、寄存器以及这些部分如何连接。Vivado再把描述转成能放进FPGA的电路结构。可综合的代码决定硬件；测试代码负责在电脑里制造输入、检查输出。

例如，两个独立的\code{always}块通常代表两部分同时工作的硬件。读源码时从上到下看，方便理解；电路运行时，这些模块会并行工作。实验4里计数器继续计数的同时，显示模块也一直在扫描。

\pagepart{0和1、位宽与二进制}
数字电路用不同电压范围表示0和1，分别称为低电平和高电平。0和1是逻辑值，高电平对应的实际电压由器件和接口标准决定；不必把所有数字电路的1都理解成同一个电压。

一位二进制数据叫一位，英文是bit。一个信号可以只有一位，例如\code{en}；也可以有多位，例如\code{input_a[3:0]}。后者有四位，从高位到低位依次是第3、2、1、0位。
\[
(0101)_2=0\times8+1\times4+0\times2+1\times1=5.
\]
四位无符号数有$2^4=16$种组合，范围是0至15。八位有256种组合，范围是0至255。这里的“无符号”表示不为负数单独分配含义。

\begin{center}
\begin{tabularx}{\linewidth}{@{}lllX@{}}\toprule
写法 & 位数 & 数值 & 读法\\\midrule
\code{1'b1} & 1 & 1 & 一位、二进制的1\\
\code{4'b0101} & 4 & 5 & 四位、二进制的0101\\
\code{4'd5} & 4 & 5 & 四位、十进制的5\\
\code{8'hA5} & 8 & 165 & 八位、十六进制的A5\\\bottomrule
\end{tabularx}
\end{center}

十六进制用0至9及A至F表示16个数字。一个十六进制数字正好对应四个二进制位，所以\code{8'hA5}也就是\code{8'b10100101}。波形中用十六进制显示八位LED图案，会比写八个二进制字符紧凑。

数字的含义取决于用途。\code{8'b00000101}若作为无符号整数，表示5；若接八个LED，则表示第0和第2个灯亮。灯图案里的每一位直接控制一盏灯，这时没必要把它转换成十进制。

\focus{位宽会限制结果。4位输出只能保留4位信息：$15+1=16=(10000)_2$，只保留低4位就得到0000。$2-5=-3$的低4位是1101；按无符号方式看是13，按4位补码方式看才是$-3$。你的实验1没有额外的符号或溢出指示灯。}

\subsection{与、或、非和选择}
“与”要求两个条件同时成立，例如允许计数且已经到上限；“或”表示至少一个条件成立；“非”把0、1反转。Verilog中的\code{&}、\code{|}、\code{~}是按位操作。对于一位信号，可以直接用这三个含义理解。

\code{sel ? value_a : value_b}表示：\code{sel=1}时选\code{value_a}，为0时选\code{value_b}。它会描述一个选择电路，后面几个实验都会用到。

\pagepart{时钟、状态和复位}
有些电路只根据当前输入给出结果，例如加法器。另一些电路要保存过去的信息，例如已经按了几次、当前亮到第几盏灯。保存的信息称为\textbf{状态}。寄存器是保存状态的一类硬件。

时钟\code{clk}周期性地从0变成1，再回到0。0变成1的瞬间叫上升沿，1变成0的瞬间叫下降沿。时序电路在约定的边沿更新寄存器，你的几个实验统一使用上升沿。

\begin{center}
\begin{tikzpicture}[x=0.65cm,y=0.75cm]
\draw[->,muted] (0,-0.4)--(18,-0.4) node[right,font=\small]{时间};
\draw[thick,ink] (0,0)--(1,0)--(1,1)--(3,1)--(3,0)--(5,0)--(5,1)--(7,1)--(7,0)--(9,0)--(9,1)--(11,1)--(11,0)--(13,0)--(13,1)--(15,1)--(15,0)--(17,0);
\foreach \x in {1,5,9,13}{\draw[flow,accent] (\x,1.8)--(\x,1.1);}
\node[anchor=west,font=\small] at (0,2.2) {箭头标记上升沿；两个相邻上升沿之间是一拍。};
\draw[<->] (5,-0.9)--(9,-0.9) node[midway,below,font=\small]{一个周期$T$};
\end{tikzpicture}
\end{center}

频率$f$表示每秒有多少个周期，周期$T=1/f$。100MHz等于$100\times10^6$Hz，因此
\[
T=\frac{1}{100\times10^6}\ \mathrm{s}=10\ \mathrm{ns}.
\]
1秒等于1000毫秒，1毫秒等于1000微秒，1微秒等于1000纳秒。你按住按钮一秒，内部时钟已经走过一亿拍。

时钟沿之间，寄存器的输出保持；它的输入可以由组合电路不断计算。到了下一个时钟沿，寄存器保存输入值。一个常见结构就是“用旧状态计算下一状态，再在时钟沿一起更新”。

\subsection{复位与使能承担不同职责}
复位把状态放回指定起点。你的\code{rst}和\code{clr}都按高电平有效设计：为1时复位。异步复位意味着复位一旦有效，寄存器就响应，不需要等下一个时钟沿。实际硬件仍有传播延迟。

使能决定这一拍是否允许更新。例如\code{en=0}时，D触发器保持原值；\code{cnt_end=0}时，流水灯保持位置。禁用后时钟仍然在运行。

组合电路没有要恢复的历史状态，因此实验1不需要复位。时序电路需要从已知状态开始，实验2清零存储内容，实验3把灯放回LED0，实验4把各计数清零并恢复自动运行。
真实寄存器要求输入在时钟沿之前和之后分别保持一小段时间，对应建立时间和保持时间。输入变化若发生在采样沿附近，就可能影响寄存器稳定地采到0或1。
